% Capa, folha de rosto e página de créditos
\begin{titlepage}
\thispagestyle{vazia}
\begin{tikzpicture}[remember picture, overlay]
  \fill[tinta] (current page.south west) rectangle (current page.north east);
  \node[anchor=south west, inner sep=0pt] at ([yshift=18mm]current page.south west)
    {\includegraphics[width=\paperwidth]{capa_leque.pdf}};
  \fill[tinta] (current page.south west) rectangle ([yshift=34mm]current page.south east);
  \node[anchor=north west, inner sep=0pt] at ([xshift=22mm, yshift=-24mm]current page.north west)
    {\includegraphics[height=17mm]{logoLiga_branco.png}};
  \node[anchor=north west, inner sep=0pt, text=white, align=left, text width=160mm]
    at ([xshift=22mm, yshift=-72mm]current page.north west)
    {{\interdisplay\fontsize{52}{55}\selectfont\espacotitulo\raggedright Finanças\\quantitativas\\com Python\par}\vspace{8mm}%
     {\interdisplay\itshape\fontsize{19}{24}\selectfont\espacotitulo\color{lilasmedio}\raggedright
      Do dado à estratégia: estatística, risco,\\backtesting e machine learning\par}};
  \node[anchor=south west, inner sep=0pt, text=white, align=left, text width=120mm]
    at ([xshift=22mm, yshift=15mm]current page.south west)
    {{\sffamily\bfseries\fontsize{13}{16}\selectfont João Emanuel Zilli dos Santos\par}\vspace{1.5mm}%
     {\sffamily\fontsize{9.5}{13}\selectfont\color{lilasmedio}Núcleo de Análise Quantitativa\par}};
  \node[anchor=south east, inner sep=0pt, text=lilasmedio]
    at ([xshift=-22mm, yshift=15mm]current page.south east)
    {\sffamily\fontsize{9.5}{13}\selectfont 1ª edição, 2026};
\end{tikzpicture}
\null
\end{titlepage}

% ------------------------------------------------------------ folha de rosto
\clearpage
\thispagestyle{vazia}
\begingroup
\vspace*{22mm}
\noindent{\interdisplay\fontsize{30}{34}\selectfont\espacotitulo Finanças quantitativas\\[1mm]com Python\par}
\vspace{5mm}
\noindent{\sffamily\large\color{cinza}Do dado à estratégia: estatística, risco,\\backtesting e machine learning\par}
\vspace{26mm}
\noindent{\sffamily\bfseries\Large João Emanuel Zilli dos Santos\par}
\vspace{2mm}
\noindent{\sffamily\small Estudante de Engenharia de Computação com IA Aplicada\par}
\vfill
\noindent\includegraphics[height=15mm]{logoLiga_preto.png}\par
\vspace{5mm}
\noindent{\sffamily\small Material didático do Núcleo de Análise Quantitativa\\
Liga de Mercado Financeiro, Universidade Federal de Uberlândia\\[2mm]
1ª edição, 2026\par}
\endgroup
\clearpage

% ------------------------------------------------------------ créditos
\thispagestyle{vazia}
\begingroup\small\setlength{\parindent}{0pt}\setlength{\parskip}{2.5mm}
\vspace*{\fill}
{\sffamily\bfseries Finanças quantitativas com Python}\\
João Emanuel Zilli dos Santos

Material didático do Núcleo de Análise Quantitativa da Liga de Mercado Financeiro da Universidade Federal de Uberlândia (LMF~UFU), escrito a partir do ciclo de aulas de 2026.1 e da ementa proposta para o Núcleo. 1ª edição, 2026.

\textbf{Aviso.} Este livro tem finalidade exclusivamente educacional. Nada nele é recomendação de investimento. As estratégias aparecem como exemplos de método, e os resultados históricos mostrados não garantem resultados futuros. Operar no mercado envolve risco de perda do capital.

\textbf{Dados.} Os preços das ações e do fundo BOVA11 vêm do histórico oficial de cotações da B3 (arquivos COTAHIST), de janeiro de 2010 a dezembro de 2025, corrigidos por desdobramentos, grupamentos, bonificações e uma cisão (a da XPart, em ITUB4) segundo comunicados das companhias. O ajuste por proventos usa o fator distribuído pelo Yahoo Finance, com as correções conferidas nesses comunicados. O Ibovespa, o S\&P~500 e o dólar comercial vêm do Yahoo Finance, e as séries do CDI e do dólar PTAX, do Sistema Gerenciador de Séries Temporais do Banco Central do Brasil. A cópia bruta de cada fonte e as tabelas de eventos e de correções estão em \arq{livro/dados/bruto/}; os scripts \arq{extrair\_cotahist.py}, \arq{limpar\_dados.py}, \arq{validar\_dados.py} e \arq{conferir\_fonte.py}, em \arq{livro/scripts/}, refazem e conferem os arquivos usados nos capítulos, como descreve o Capítulo~\ref{cap:dados}. Os dados foram obtidos em outubro de 2026 e têm uso educacional, sem garantia de exatidão.

\textbf{Repositório.} O livro, os códigos, os dados e o material das aulas estão em \url{https://github.com/joaoemanuelzilli/aulas_quant}. Os caminhos de arquivo citados no texto, como \arq{aula5/Var\_historico.py} ou \arq{livro/scripts/cap05/}, são relativos à raiz desse repositório.

\textbf{Código.} Os códigos impressos e as figuras de dados são gerados por scripts em \arq{livro/scripts/}, uma pasta por capítulo, executados com Python~3.12 e as versões de bibliotecas registradas em \arq{livro/requirements.txt}. Cada script roda de dentro da própria pasta (\cod{cd livro/scripts/cap01} e depois \cod{python l02\_pandas.py}), porque lê os dados por um caminho relativo a ela. Os poucos trechos apresentados como esboço, que dependem de serviços externos, estão identificados no texto.

\textbf{Composição.} Texto em STIX Two, títulos em Inter e código em JetBrains Mono, fontes distribuídas sob a SIL Open Font License. Composto em Lua\LaTeX.

A marca e o logotipo da Liga de Mercado Financeiro pertencem à Liga e são usados neste material institucional com essa finalidade.
\endgroup
\clearpage
